% Generated draft: requires mathematical and source review.
\documentclass[11pt]{article}
\usepackage[margin=25mm]{geometry}
\usepackage{fontspec}
\setmainfont{Latin Modern Roman}
\usepackage{amsmath,amssymb,mathtools,unicode-math}
\setmathfont{Latin Modern Math}
\usepackage{enumitem,needspace,xurl,hyperref}
\setlength{\parindent}{0pt}
\setlength{\parskip}{5pt}
\setlength{\emergencystretch}{3em}
\newcommand{\N}{\mathbb{N}}
\newcommand{\Z}{\mathbb{Z}}
\newcommand{\Q}{\mathbb{Q}}
\newcommand{\R}{\mathbb{R}}
\newcommand{\C}{\mathbb{C}}
\newcommand{\F}{\mathbb{F}}
\newcommand{\A}{\mathbb{A}}
\newcommand{\PP}{\mathbb{P}}
\newcommand{\E}{\mathbb{E}}
\newcommand{\eps}{\varepsilon}
\newcommand{\dd}{\,\mathrm{d}}
\begin{document}
\section*{UnsolvedMath ID: 6000010}
\textbf{Source problem code:} catalogue entry
\subsection{Definitions and mathematical statement}
Let $(X^{2n},\omega,J)$ be a symplectic manifold with compatible almost complex structure: $J^2=-1$, $\omega(JU,JV)=\omega(U,V)$, and $\omega(U,JU)>0$ for $U\ne0$. No integrability of $J$ is assumed.

For a Riemannian manifold $(M,g)$, the metric identification $b_g:TM\to T^*M$, $v\mapsto g(v,\cdot)$, pulls the canonical cotangent symplectic form back to $\omega_g=b_g^*\omega_{\mathrm{can}}$. A torsion-free affine connection $\nabla$ splits $T(TM)$ into horizontal and vertical copies of $TM$. In this splitting its standard Dombrowski-type almost complex structure is
\[
 J_\nabla(U,V)=(-V,U).
\]
The source also requires the connection to be ``compatible,'' but does not say whether this means metric compatibility, a statistical/Codazzi condition, or precisely compatibility of $J_\nabla$ with $\omega_g$. Those interpretations are not silently equated.

\textbf{Characterization problem (Takashi Kurose).} Give intrinsic necessary and sufficient conditions for $(X,\omega,J)$ to be locally of the above tangent-bundle form. More explicitly, near each relevant point seek a local diffeomorphism $\Phi$ into an open subset of $TM$, for suitable $M,g,\nabla$, satisfying
\[
 \Phi^*\omega_g=\omega,\qquad
 d\Phi\circ J=J_\nabla\circ d\Phi,
\]
with the intended compatibility condition on the torsion-free connection. The desired result is an intrinsic criterion/classification, not merely this restatement of local realization. Whether local charts must meet a distinguished zero section or preserve further tangent-bundle data is unspecified.

The supplied summary describes the construction and classical structure equations as understood, naming Dombrowski-type geometry, but says no fully general classification beyond those equations was verified. It gives no exact theorem or definition of compatible connection. This is attributed background, not certification that the characterization requested here is complete.
\subsection{Short English statement}
Characterize compatible symplectic/almost-complex structures locally modeled on a metric tangent bundle with the connection-induced Dombrowski structure. The meaning of connection compatibility and the required local tangent-bundle data are unspecified. Almost complex is not strengthened to complex or Kähler.
\subsection{Sources}
\begin{itemize}
\item Local UnsolvedMath registry, permanent ID 6000010, source code catalogue entry, ``Tangent-Bundle Symplectic and Almost Complex Structures''; proposer: Takashi Kurose, proposed year 1998. Status open, classification OPEN-TRIAGE, updated 2026-08-06.
\item Furuhata, Matsuzoe, and Urakawa, \emph{Open Problems in Affine Differential Geometry and Related Topics} (1998), item 3(d), scanned PDF pages 1--2, \url{https://www.jstage.jst.go.jp/article/iis/4/2/4_2_125/_pdf/-char/en}. Reported manual OCR from page image and access: 2026-07-29. The article presents an open item; current status NEEDS\_REVIEW. The record says the J-STAGE article page marks it CC BY 4.0.
\item The summary names Dombrowski-type almost complex structures and classical tangent-bundle structure equations without identifying a publication or full characterization theorem.
\end{itemize}
\end{document}
